\documentclass[11pt,a4paper]{article}
\usepackage[utf8]{inputenc}
\usepackage[T1]{fontenc}
\IfFileExists{polish.ldf}{\usepackage[polish]{babel}}{}
\IfFileExists{lmodern.sty}{\usepackage{lmodern}}{}
\usepackage[margin=26mm]{geometry}
\usepackage{amsmath}
\usepackage[authoryear,round]{natbib}
\usepackage[hidelinks]{hyperref}
\title{Kandydat: publikacja metod oprogramowania EEG Neurolyze}
\author{Mateusz Pomianek}
\date{MSC-EEG-P01 | Manuskrypt metod 1.0.0 | 27 listopada 2026}
\begin{document}
\maketitle
\begin{abstract}
EEG Neurolyze wspiera demonstrację analizy widmowej różnic między spoczynkiem z oczami otwartymi i zamkniętymi. Przedstawiono wersjonowany przebieg od wygenerowania i archiwizacji sygnału po odtwarzalną analizę mocy alfa. W 18 syntetycznych przypadkach porównano implementację Welcha z SciPy przy wyrównanych parametrach. Największa bezwzględna różnica gęstości widmowej wyniosła $3{,}98\times10^{-13}\,\mu\mathrm{V}^2/\mathrm{Hz}$. Dowody dotyczą zdefiniowanych wejść programowych.
\end{abstract}
\section{Wprowadzenie}
Reprodukowalność EEG obejmuje parametry akwizycji, przetwarzania, opis metadanych i dostępność specyfikacji analizy \citep{pernet2020}. EEG-BIDS porządkuje reprezentację danych, zdarzeń i metadanych \citep{pernet2019}. Wkład Neurolyze rozpatrujemy jako spójność lokalnego przebiegu dydaktycznego i możliwość prześledzenia wyniku od próbek do raportu. Odrębnym punktem odniesienia jest MNE-BIDS-Pipeline, który dostarcza ogólną automatyzację przetwarzania zbiorów BIDS. Artykuł MSC-EEG-P02 obejmuje kwalifikację urządzeń, w tym pytania o synchronizację, mapowanie kanałów i stabilność pomiaru.

\section{Materiał i metody}
Oprogramowanie generuje próbki syntetyczne oraz oznaczenia warunków EO i EC, zapisuje dane, informacje o kanałach, zdarzenia i manifest SHA-256 w archiwum EEGZ oraz pozwala obliczyć analizę ponownie z zapisu. Sumy kontrolne służą sprawdzeniu integralności. W scenariuszu przykładowym po stabilizacji (30 s) następuje EO (120 s), przerwa (30 s) i EC (120 s). Skrócone sesje wykorzystano w testach integracyjnych. Eksport BrainVision EEG-BIDS obejmuje wejścia syntetyczne z równomiernym próbkowaniem. Zgodność BIDS wymaga wyniku walidatora dla konkretnego zbioru.

Gęstość widmowa jest estymowana metodą Welcha \citep{welch}. Algorytm usuwa średnią z każdego segmentu, stosuje symetryczne okno Hanna i krok równy połowie długości segmentu po zaokrągleniu w dół. Segmenty mają nominalnie 2 s, przy krótszym wejściu używana jest jego długość. FFT jest dopełniana zerami do najbliższej potęgi dwóch. Widmo jednostronne ma jednostkę $\mu\mathrm{V}^2/\mathrm{Hz}$, a moc alfa liczymy jako
\[
 P_\alpha=\Delta f\sum_{k:\,8\leq f_k\leq 13} \widehat S(f_k),
\]
gdzie $\Delta f$ oznacza odstęp między częstotliwościami siatki FFT. Różnice w wyborze okna, nakładaniu, obsłudze braków oraz granicach pasma mogą zmieniać wynik.

\section{Plan analizy}
Pierwszorzędowy wskaźnik: zgodność wartości PSD i mocy alfa między Neurolyze a SciPy przy wyrównanych ustawieniach. Jednostką porównania jest syntetyczny sygnał o zadanej częstotliwości próbkowania i długości. Zestaw skrzyżował 128, 200 i 256 Hz z dwoma długościami (około 1 i 5 s) oraz trzema typami sygnału. Były to sinus 10 Hz o amplitudzie 20 $\mu$V, sinus 10 Hz z dryfem i szumem oraz sygnał zerowy. SciPy otrzymało jawne okno symetryczne i dokładnie dopasowane parametry. Sprawdzono identyczność siatek oraz zgodność PSD i całki pasmowej z tolerancjami bezwzględną oraz względną $10^{-10}$. Zamrożenie estymandu dla analizy empirycznej pozostaje odrębną decyzją po dopuszczeniu zbioru.

\section{Wyniki}
Wszystkie 18 przypadków spełniło kryterium zgodności. Maksimum bezwzględnej różnicy PSD wyniosło $3{,}98\times10^{-13}\,\mu\mathrm{V}^2/\mathrm{Hz}$. Dla sinusa moc 8--13 Hz mieściła się w przedziale 199,935--200,000 $\mu\mathrm{V}^2$ po zaokrągleniu do trzech miejsc. Przeszło także dziewięć testów widmowych i trzy testy integracyjne EO/EC przypiętego źródła. Wśród nich są przypadki wywołania CLI, archiwizacji i błędów znaczników.

\section{Dyskusja i ograniczenia}
Zgodność z referencją wspiera poprawność numeryczną dla badanych długości, częstotliwości i typów wejścia. Benchmark dotyczy równomiernego próbkowania. Odczyt publicznego EDF w oddzielnym przykładzie MNE, kwalifikacja Muse S Athena i walidacja konkretnego zbioru BIDS wymagają osobnych dowodów. Przed zgłoszeniem należy przypiąć wydanie oprogramowania i udokumentować instalację, licencję oraz odtworzenie procedury przez niezależnego użytkownika. Docelowy format to artykuł o oprogramowaniu badawczym.
\bibliographystyle{apalike}
\bibliography{references}
\end{document}
